% Part I, ends with ICFP 2016 (no divider: part of program calculation)

\begin{frame*}[vertical-alignment = top]
  \frametitle{One Specification, Many Parallel Programs}
  \vspace*{3mm}
  \hsblock[before upper = {\fvset{fontsize = \footnotesize}}]{psort.hs}
  \begin{center}
    \begin{tikzpicture}[alt={A three-stage parallel sort. A list enters a
        split stage, which halves every list longer than n and feeds the
        halves back into its own input. The short lists go to a farm of k
        workers that sort them. The sorted lists go to a merge stage, which
        merges two sorted lists and feeds the result back, until one sorted
        list remains.},
      stage/.style = {box, font = \small, minimum height = 7mm,
        minimum width = 12mm, inner sep = 3pt},
      loop/.style = {->, accent2, thick},
      >={Stealth[length = 1.5mm]}]
      \node (in) {list};
      \node[stage, right = 8mm of in] (sp) {split};
      \foreach \i/\y in {1/5, 2/0, 3/-5}
        \node[stage, right = 16mm of sp, yshift = \y mm, minimum height = 5mm]
          (w\i) {sort};
      \node[stage, right = 16mm of w2] (mg) {merge};
      \node[right = 8mm of mg] (out) {sorted};
      \draw[->] (in) -- (sp);
      \foreach \i in {1, 2, 3} { \draw[->] (sp) -- (w\i);
        \draw[->] (w\i) -- (mg); }
      \draw[->] (mg) -- (out);
      \draw[loop] (sp.north east) .. controls +(0.6, 0.8) and +(-0.6, 0.8)
        .. (sp.north west);
      \draw[loop] (mg.north east) .. controls +(0.6, 0.8) and +(-0.6, 0.8)
        .. (mg.north west);
      \node[note, above = 6mm of sp] {halves fed back};
      \node[note, above = 6mm of mg] {merged pairs fed back};
      \node[note, above = 0mm of w1] {farm of $k$};
    \end{tikzpicture}
  \end{center}
  \begin{onlyenv}<2->
    {\centering $\op{psort} : \List\,\op{Int} \xrightarrow{\;\sigma\;}
      \List\,\op{Int}$ \qquad $\sigma = \op{FB}\,(\op{FUN}\,A) \parallel
      \op{FARM}\,k\,(\op{FUN}\,A) \parallel \op{FB}\,(\op{FUN}\,A)$\par}
    \smallskip
    The type carries the structure drawn above; type checking proves
    \code{psort} equal to \code{msort}, with the laws. A \textbf{cost
    model} searches for the \textbf{fastest} $\sigma$, $n$ and $k$.
  \end{onlyenv}
  \slideref{D. Castro, K. Hammond, S. Sarkar. \textbf{Farms, Pipes, Streams and Reforestation: Reasoning about Structured Parallel Processes using Types and Hylomorphisms}. ICFP 2016. \rfsep
    D. Castro, K. Hammond, S. Sarkar, Y. Alguwaifli. \textbf{Automatically Deriving Cost Models for Structured Parallel Processes Using Hylomorphisms}. FGCS 2018.}
  \icfptag
  \speakernote{Skeletons: reusable parallel patterns; fb = feedback, farm
    = replicated workers, || = pipeline \\
    Split stage: an unfold as a loop; merge stage: a fold as a loop \\
    Same result as mergesort, many choices of n and k \\
    sigma: FB = feedback, FARM k = k workers, || = pipeline, FUN =
    sequential \\
    Is it equal? The type checker answers, normalising with the laws
    (reforestation: put the tree back to expose stages) \\
    Which is fastest? A cost model, from the same structure (FGCS 2018) \\
    Other structures (splitter and merger farms), the skeleton table,
    reforestation: backup}
\end{frame*}

\begin{frame}[vertical-alignment = top]
  \frametitle{What Was Missing}
  \vspace*{2mm}
  \begin{itemize}[<+->]
    \item \textbf{Laws as axioms}: the ICFP 2016 type system assumed them,
      and programs could diverge.
    \item \textbf{Side conditions on trust}: \code{hylo merge halves}
      becomes \code{... || fb (fun merge)} only if \code{merge} is
      associative and commutative; the library took the programmer's word
      for it.
    \item \textbf{Unverified tools}: implement the patterns or rewrites,
      then calculate by hand or by trusted rewriting.
    \item \textbf{No compilation to idiomatic code}: calculated programs
      stay higher-order; we trust the compiler to inline them into
      efficient code.
  \end{itemize}
  \uncover<5->{\begin{center}
    \begin{tikzpicture}[alt={The question for the rest of the talk: can a
        proof assistant prove each law once, check every rewrite, side
        conditions included, and still produce efficient code?}]
      \node[hot, text width = 140mm, align = center, inner sep = 1.5mm]
        {Can a proof assistant \textbf{prove} each law once,
        \textbf{check} every rewrite, side conditions included, and still
        produce \textbf{efficient} code?};
    \end{tikzpicture}
  \end{center}}
  \slideref{Tools: HYLO (Hu et al., ICFP'96); MAG (de Moor, Sittampalam, TCS'01); Pointless Haskell, DrHylo (Cunha et al., IFL'05); SkeTo (Matsuzaki, Emoto, HLPP'10); HERMIT (Farmer et al., Haskell'12).}
  \icfptag
  \speakernote{ICFP'16: soundness on paper, laws taken for granted; cpo
    semantics \\
    Side conditions: merging in a feedback loop pairs chunks in any order,
    so merge must be associative and commutative; never checked: the
    PyTorch failure mode \\
    Same elsewhere: compilers (HYLO), transformation systems (MAG, HERMIT),
    point-free tools (Cunha et al.), skeleton libraries (SkeTo) \\
    Higher-order code (folds, combinators, skeletons): efficiency left to
    the compiler's inliner and code generator, both trusted \\
    Next: can a proof assistant do this?}
\end{frame}
